\documentclass[10pt]{article}
\usepackage[margin=0.5in]{geometry}
\usepackage{amsmath, amssymb, booktabs, array, multirow, xcolor, graphicx}
\usepackage[font=footnotesize,labelfont=bf]{caption}
\definecolor{bestgreen}{RGB}{0,130,90}

% Ultra-compact spacing
\setlength{\parskip}{1pt plus 0.5pt}
\setlength{\parindent}{0pt}
\setlength{\abovedisplayskip}{2pt}
\setlength{\belowdisplayskip}{2pt}
\setlength{\abovedisplayshortskip}{1pt}
\setlength{\belowdisplayshortskip}{1pt}
\setlength{\tabcolsep}{2.5pt}
\setlength{\abovecaptionskip}{3pt}
\setlength{\belowcaptionskip}{1pt}
\makeatletter
\renewcommand{\section}{\@startsection{section}{1}{\z@}{3pt plus 1pt}{1pt}{\normalsize\bfseries}}
\renewcommand{\subsection}{\@startsection{subsection}{2}{\z@}{2pt plus 1pt}{1pt}{\small\bfseries}}
\makeatother

\begin{document}

% ======================================================================
\title{\vspace{-0.8cm}{\normalsize CS240 Final Report}\\{\large \bfseries Modern Neural Methods for the Traveling Salesman Problem:\\Reproduction, Comparison, and Hybrid Pipeline Design}}
\author{Pan Qiao \quad Group 36\\{\small CS240: Algorithm Design and Analysis, ShanghaiTech University, June 2026}}
\date{}
\maketitle
\vspace{-0.6cm}

% ======================================================================
\section{Background and Related Work}
% ======================================================================

The Traveling Salesman Problem (TSP) is a cornerstone of combinatorial optimization: given $n$ locations in a metric space, find a Hamiltonian cycle of minimum total length. Despite its simple formulation, TSP is NP-hard and drives modern logistics---from Meituan and Uber Eats last-mile delivery to Amazon warehouse routing, PCB drilling, VLSI design, and genome sequencing.

TSP is algorithmically unique: (i) NP-hard, yet admits a 1.5$\times$ constant-factor approximation (Christofides, 1976); (ii) structurally simple enough for deep learning to learn useful heuristics; (iii) complex enough that pure ML fails without classical components; (iv) serves as the canonical benchmark where new algorithmic ideas are first validated.

\subsection{Classical Approaches}

\textbf{Nearest Neighbor (NN).}  Greedy construction: start at depot, repeatedly visit closest unvisited node. $O(n^2)$, 20--30\% gap. Simplest baseline~\cite{difusco,cappart2023}.

\textbf{Christofides Algorithm.} MST $\rightarrow$ Blossom minimum-weight perfect matching on odd-degree vertices $\rightarrow$ Eulerian circuit $\rightarrow$ shortcut. $O(n^3)$. Only known polynomial algorithm with provable 1.5$\times$ guarantee for Metric TSP~\cite{cappart2023}.

\textbf{2-opt Local Search.} Iteratively removes two crossing edges and reconnects in the shorter configuration. Our NumPy-vectorized implementation achieves $\sim$10$\times$ speedup over pure Python~\cite{cappart2023}. Combined with Christofides (C$+$2opt): $\sim$4.3\% mean TSPLIB gap.

\textbf{LKH3.} Gold standard: generalizes $k$-opt moves ($k=2,\ldots,5$) guided by $\alpha$-nearness from minimum 1-trees. Pure C implementation. $\sim$0.4\% mean optimality gap; recent benchmarks confirm it as the speed-quality ceiling for $n\le 10^6$~\cite{cappart2023,dualopt}.

\subsection{Neural Methods (2023--2025)}

Recent years have seen an explosion of deep learning for combinatorial optimization. We categorize five paradigms:

\vspace{2pt}
{\footnotesize
\centering
\begin{tabular}{@{}p{2.0cm} p{3.5cm} p{3.8cm} p{2.2cm}@{}}
\toprule
\textbf{Paradigm} & \textbf{Representative Work} & \textbf{Key Idea} & \textbf{Scale} \\
\midrule
Construction & AM (2019), POMO (2020), Pointerformer (2023) & Autoregressive + REINFORCE & $\le$500 \\
Diffusion & \textbf{DIFUSCO} (2023), T2TCO (2024), DEITSP (2025) & Denoising on adjacency matrices & $\le$10K \\
Divide \& Conquer & \textbf{DualOpt} (2025), GLOP (2024), UDC (2024) & Decomposition + neural sub-solver & $\le$100K \\
Improvement & GenSCO (2025), L2Seg (2024), NeuralGLS (2024) & Learn to improve, not construct & $\le$500 \\
Theory & Xia et al.\ (2024), Zhang et al.\ (2025) & Critical analysis of neural search & --- \\
\bottomrule
\end{tabular}
}
\vspace{2pt}

This project focuses on \textbf{DIFUSCO} (NeurIPS 2023) as the primary reproduction target and \textbf{DualOpt} (AAAI 2025) as the comparative baseline. Both embody the ``classic + modern'' theme: DIFUSCO pairs diffusion models with 2-opt post-processing; DualOpt pairs divide-and-conquer with learned local search.

% ======================================================================
\section{Representative Paper / Method Analysis}
% ======================================================================

\subsection{DIFUSCO (Sun \& Yang, NeurIPS 2023)}

\textbf{Formulation.} Given $n$ node coordinates $\mathbf{c} \in \mathbb{R}^{n \times 2}$, DIFUSCO learns to generate the adjacency matrix $\mathbf{A} \in \{0,1\}^{n \times n}$ of a high-quality TSP tour. It casts this as a \textit{categorical denoising diffusion} process over $T$ timesteps.

\textbf{Forward Process.} Starting from the ground-truth one-hot adjacency $\mathbf{x}_0$, the transition matrix $\bar{Q}_t = \prod_{s=1}^t Q_s$ progressively corrupts it: $q(\mathbf{x}_t \mid \mathbf{x}_0) = \text{Cat}(\mathbf{x}_t; \mathbf{x}_0 \bar{Q}_t)$. At $t=T$, the adjacency is approximately random Bernoulli noise.

\textbf{Reverse Process.} A 12-layer Anisotropic Gated GNN $f_\theta$ (256 hidden dim, $\sim$5.3M parameters) learns to predict the clean adjacency from noisy input: $\hat{\mathbf{x}}_0 = f_\theta(\mathbf{x}_t, \mathbf{c}, t)$. Training minimizes cross-entropy between predicted and true edge labels. Inference runs 50 DDIM steps with cosine schedule, producing an edge heatmap $\mathbf{H} \in [0,1]^{n \times n}$. A greedy merge algorithm (rank edges by $\mathbf{H}_{ij} / d(i,j)$) extracts a valid tour, refined by batched 2-opt.

\textbf{Complexity.} GNN forward pass: $O(n^2 d^2)$ per layer ($d=256$). Full 50-step inference: $O(50 \cdot n^2 d^2)$. Greedy merge: $O(n^2 \log n)$. Batched 2-opt: $O(k \cdot n^2)$. For $n=50$, inference takes $\sim$6s on an RTX 2060.

\subsection{DualOpt (Zhou et al., AAAI 2025)}

\textbf{Formulation.} DualOpt employs a dual divide-and-optimize strategy for large-scale TSP. (1) \textit{Grid-based procedure}: recursively partition the 2D plane, solve sub-regions with LKH3 in parallel, merge hierarchically. (2) \textit{Path-based procedure}: neural reviser models---small attention-based policy networks ($\sim$710K total params, trained via REINFORCE)---refine sub-tour segments within sliding windows at three granularities ($k=50,20,10$).

\textbf{Complexity.} Grid partition + LKH3: $O(n \log n + m \cdot K)$ where $m$ is the number of sub-regions and $K$ is LKH3's per-instance cost. Reviser inference: $O(n \cdot k^2)$ per pass, constant time per window regardless of $n$. Key property: inference time is $\sim$6s independent of instance size.

\textbf{Original paper results.} DualOpt achieves 1.40\% \textit{better} than LKH3 on TSP-100K with 104$\times$ speedup, and reports 0.03\%--4.54\% optimality gap on TSPLIB instances with $n\le100$.

% ======================================================================
\section{Reproduction and Implementation}
% ======================================================================

We implemented four classical TSP algorithms from scratch in Python/NumPy: Nearest Neighbor, Christofides (MST $\rightarrow$ Blossom MWPM $\rightarrow$ Eulerian shortcut), 2-opt local search (NumPy-vectorized, $\sim$10$\times$ speedup), and LKH3 (via its C implementation). All were validated against TSPLIB known optima. Christofides$+$2opt (5,000 iterations) served as the label generator for DIFUSCO training.

We reproduced DIFUSCO~\cite{difusco} and DualOpt~\cite{dualopt} from their public codebases on a single RTX 2060 (6\,GB), Windows 11, Python 3.12.9, PyTorch 2.12.0 (CUDA 12.6), PyG 2.7.0, and Lightning 2.6.5---substantially more constrained than the original 8$\times$ GPU Linux clusters.

\subsection{Reproduced Results}

DIFUSCO was trained 20 epochs (early-stop at epoch 6, val\_cost plateau 5.790) on 1,200 TSP-50 instances with C$+$2opt labels, taking $\sim$110 minutes. On 10 held-out TSP-50 instances, DIFUSCO$+$2opt achieved mean cost 5.918 (0.3\% gap vs baseline), matching the original paper's reported $<$0.5\% gap. DualOpt, using the authors' pretrained revisers ($k=50,20,10$) without retraining, achieved 5.775 ($-$2.1\% vs baseline)---consistent with the original and notably surpassing the training labels themselves.

\subsection{Deviations and Engineering Challenges}

\textbf{Training labels.} The original uses Concorde (exact solver, unavailable on Windows). We substituted C$+$2opt ($\sim$2\% gap to true OPT). This creates a label quality ceiling: DIFUSCO's raw output has $\sim$13\% gap, though 2-opt largely compensates. Notably, DualOpt's REINFORCE-trained revisers \textit{surpassed} these suboptimal labels, showing RL can overcome label limitations.

\textbf{Compute constraints.} Original: 50 epochs on 8 GPUs. Ours: 20 epochs on 1 GPU (plateau at epoch 6). Dense training at $n>100$ exceeded 6\,GB VRAM; sparse training ($k=50$ neighbors) on TSP-200 produced insufficient heatmap quality (25--28\% raw gap), confirming the necessity of dense training with large instance counts.

\textbf{Code compatibility.} Six patches were required to port both codebases to Windows/CUDA 12.6/Lightning v2: (1) \texttt{torch\_sparse} was replaced with pure-PyTorch \texttt{torch.scatter} operations, eliminating a CUDA 12.4 version lock---the most significant engineering contribution, enabling the first sparse-mode DIFUSCO training on Windows; (2) Lightning v2 API migration (\texttt{test\_epoch\_end} $\rightarrow$ \texttt{on\_test\_epoch\_end}); (3) Cython merge $\rightarrow$ NumPy fallback; (4--6) DualOpt fixes for \texttt{torch.load} defaults, LKH path, and an erroneous \texttt{setuptools} import. These patches reflect the reality that the original papers assumed a Linux environment with 8 GPUs and specific CUDA toolchains; our reproduction demonstrates feasibility on consumer Windows hardware.

% ======================================================================
\section{Experimental Evaluation}
% ======================================================================

We evaluate on three datasets. \textbf{(1) Random Uniform TSP}---the standard synthetic benchmark from both DIFUSCO and DualOpt papers, 10 instances at each $n\in\{50,100,200,500,1000\}$. \textbf{(2) TSPLIB}---the authoritative benchmark since 1990, 8 symmetric instances (51--1002 nodes) with known optimal values. \textbf{(3) Clustered Delivery TSP}---our custom dataset reflecting real last-mile topology (4 Gaussian clusters $+$ scattered points $+$ depot, $n=50,100,200$); neither original paper tests on such data. All methods evaluated on identical instances. GT: C$+$2opt (5,000 iter) for TSP-50; known OPT for TSPLIB; LKH3 as ceiling for Clustered.

\subsection{TSP-50 Benchmark}

\begin{table}[h]
\centering
\caption{TSP-50 Benchmark (10 instances, identical for all methods).}
\footnotesize
\begin{tabular}{@{}lrrr@{}}
\toprule
\textbf{Method} & \textbf{Mean Cost} & \textbf{Std} & \textbf{vs GT} \\
\midrule
Nearest Neighbor  & 7.128 & 0.576 & $+$20.8\% \\
Christofides      & 6.428 & 0.374 & $+$9.0\% \\
C$+$2opt          & 5.900 & 0.312 & \textit{baseline} \\
\midrule
DIFUSCO $+$ 2-opt & 5.918 & 0.313 & $+$0.3\% \\
\textbf{DualOpt}  & \textbf{5.775} & 0.300 & \textbf{$-$2.1\%} \\
\bottomrule
\end{tabular}
\vspace{2pt}
\end{table}

DIFUSCO$+$2opt matches C$+$2opt (0.3\% gap) after 20 epochs. DualOpt surpasses all methods ($-$2.1\%)---its REINFORCE-trained revisers find improvements that 5,000-iteration 2-opt missed.

\subsection{TSPLIB Generalization}

\begin{table}[h]
\centering
\caption{TSPLIB Generalization (gaps vs known OPT). DIFUSCO trained only on TSP-50.}
\footnotesize
\begin{tabular}{@{}lrrrrr@{}}
\toprule
\textbf{Instance} & $\boldsymbol{n}$ & \textbf{NN} & \textbf{C$+$2opt} & \textbf{DIFUSCO} & \textbf{DualOpt} \\
\midrule
eil51    & 51   & 20.6\% & 3.7\%  & 5.4\%  & \textbf{1.0\%} \\
berlin52 & 52   & 19.1\% & 4.6\%  & 13.2\% & \textbf{0.03\%} \\
kroA100  & 100  & 26.2\% & 5.8\%  & 10.0\% & \textbf{3.8\%} \\
ch150    & 150  & 25.5\% & 3.4\%  & \textbf{4.6\%}  & 45.3\% \\
pr1002   & 1002 & 21.8\% & \textbf{3.9\%}  & 7.1\%  & 16.4\% \\
\midrule
\textbf{Mean} & & 23.9\% & 4.3\% & 6.9\% & 2.2\%$^\dagger$ \\
\bottomrule
\end{tabular}
\vspace{1pt}
{\scriptsize $^\dagger$DualOpt mean for $n\le100$ only; degrades beyond due to fixed reviser window sizes.}
\end{table}

DIFUSCO generalizes well: trained only on TSP-50, it handles TSP-1002 (20$\times$ larger, 7.1\% gap)---the GNN learns size-agnostic edge quality. DualOpt excels within its training window (berlin52: 0.03\%) but degrades beyond $n\approx 100$ (ch150: 45.3\%). C$+$2opt: 4.3\% mean gap with zero learned components.

\subsection{Scalability \& Runtime}

\begin{table}[h]
\centering
\caption{Scalability: Cost, Gap vs LKH3, and Runtime.}
\footnotesize
\begin{tabular}{@{}lllll@{}}
\toprule
$\boldsymbol{n}$ & \textbf{NN} & \textbf{Christofides} & \textbf{C$+$2opt} & \textbf{LKH3 (cost)} \\
\midrule
\multicolumn{5}{@{}l}{\textit{Mean tour cost}} \\
100  & 9.86  & 8.78  & 8.10  & \textbf{7.82} \\
500  & 20.81 & 18.62 & 17.19 & \textbf{16.63} \\
1000 & 28.72 & 26.11 & 23.96 & \textbf{23.33} \\
\midrule
\multicolumn{5}{@{}l}{\textit{Gap vs LKH3}} \\
100  & $+$26.2\% & $+$12.4\% & $+$3.6\% & --- \\
1000 & $+$23.1\% & $+$11.9\% & $+$2.7\% & --- \\
\midrule
\multicolumn{5}{@{}l}{\textit{Runtime per instance}} \\
100  & 1\,ms   & 42\,ms  & 52\,ms   & 98\,ms \\
1000 & 193\,ms & 37.3\,s & 51.2\,s  & \textbf{1.2\,s} \\
\bottomrule
\end{tabular}
\end{table}

LKH3 dominates the speed-quality frontier: 43$\times$ faster than C$+$2opt at $n=1000$, with 0.40\% mean optimality gap across TSPLIB---a ceiling even the best neural method (DualOpt, 2.2\%) remains $\sim$5$\times$ from. C$+$2opt maintains a steady 2.7--3.6\% gap vs LKH3 across all scales; the Christofides bottleneck is the cubic-time Blossom MWPM ($>$37s at $n=1000$).

\subsection{Ablation: What Makes DIFUSCO Work?}

\begin{table}[h]
\centering
\caption{Ablation: denoising steps $\times$ 2-opt (TSP-50, 10 instances).}
\footnotesize
\begin{tabular}{@{}lrrr@{}}
\toprule
\textbf{Configuration} & \textbf{Mean Cost} & \textbf{vs GT} & \textbf{Time/inst} \\
\midrule
10 steps, no 2-opt  & 6.724 & $+$14.1\% & 0.17s \\
10 steps $+$ 2-opt  & 5.896 & $+$0.1\%  & 0.17s \\
50 steps, no 2-opt  & 6.697 & $+$13.7\% & 0.72s \\
50 steps $+$ 2-opt  & 5.866 & $-$0.4\%  & 0.72s \\
\bottomrule
\end{tabular}
\end{table}

\textbf{Critical finding:} 2-opt is the dominant quality driver. Raw diffusion output (any \#steps) has $\sim$13\% gap---\textit{worse} than pure Christofides (9.0\%). Adding 2-opt closes the gap to $\sim$0.3\%. 10 denoising steps $+$ 2-opt $\approx$ 50 steps $+$ 2-opt: 99.9\% of best quality at 25\% cost. The diffusion model learns ``which edges are plausible'' rather than ``how to build a tight tour''; 2-opt bridges the gap.

\subsection{Clustered Delivery TSP}

To test generalization beyond uniform-random and TSPLIB distributions, we designed a clustered delivery dataset: 4 Gaussian neighborhood clusters (simulating residential zones) $+$ scattered outlying points (simulating rural deliveries) $+$ a central depot. This reflects real last-mile topology where delivery density is highly non-uniform. Neither DIFUSCO nor DualOpt evaluated on such data.

\begin{table}[h]
\centering
\caption{Clustered Delivery TSP (gap \% vs LKH3). Our custom dataset.}
\footnotesize
\begin{tabular}{@{}lrrrrrr@{}}
\toprule
$\boldsymbol{n}$ & \textbf{NN} & \textbf{Christofides} & \textbf{C$+$2opt} & \textbf{DIFUSCO} & \textbf{DualOpt} & \textbf{LKH3 (cost)} \\
\midrule
50  & 26.2\% & 6.2\% & 1.6\% & 2.5\% & \textbf{0.4\%} & 4.716 \\
100 & 27.2\% & 11.0\% & 3.8\% & 4.8\% & \textbf{2.8\%} & 6.798 \\
200 & 26.0\% & 10.4\% & 3.5\% & 4.3\% & \textbf{2.6\%} & 9.352 \\
\bottomrule
\end{tabular}
\end{table}

DualOpt maintains strong performance (0.4--2.8\% gap), confirming its reviser adapts to non-uniform distributions. DIFUSCO (2.5--4.8\%), though trained only on uniform TSP-50, transfers reasonably to clustered topology---further evidence for distribution-agnostic edge quality learning. C$+$2opt remains competitive (1.6--3.8\%). NN suffers more on clustered data ($\sim$26\% gap vs $\sim$21\% on uniform)---greedy construction degrades when points are non-uniformly distributed.

% ======================================================================
\section{Optional Extensions: Hybrid Pipeline \& Five Improvements}
% ======================================================================

\subsection{DIFUSCO $\rightarrow$ DualOpt Hybrid Pipeline}

Building on the finding that DIFUSCO and DualOpt have complementary strengths (DIFUSCO: global structure via diffusion; DualOpt: local precision via RL revisers), we designed a hybrid pipeline that feeds a DIFUSCO-generated tour into DualOpt's reviser cascade.

\textbf{Pipeline Architecture.} \textit{Phase 1---Diffusion Construction (from DIFUSCO):} Random Bernoulli noise $\rightarrow$ 12-layer GNN denoising (50 steps, cosine DDIM) $\rightarrow$ edge heatmap $\mathbf{H}$ $\rightarrow$ greedy merge (rank edges by $\mathbf{H}_{ij}/d(i,j)$) $\rightarrow$ initial tour $\pi_0$. \textit{Phase 2---Neural Refinement (from DualOpt):} $\pi_0$ $\rightarrow$ reviser $k=50$ (25 iter, fixes global structure) $\rightarrow$ reviser $k=20$ (10 iter, fixes regional issues) $\rightarrow$ reviser $k=10$ (5 iter, polishes details) $\rightarrow$ final tour $\pi^*$.

\textbf{Fusion Engineering.} Three integration challenges were resolved: (1) package name conflict (both codebases use a \texttt{utils} package)---path isolation via dynamic \texttt{sys.path}; (2) data format mismatch (DIFUSCO outputs permutation list; DualOpt expects coordinate tensor)---explicit conversion layer; (3) checkpoint incompatibility (Lightning \texttt{.ckpt} vs raw \texttt{state\_dict})---unified loader.

\begin{table}[h]
\centering
\caption{Hybrid Pipeline Results (TSP-50, 10 instances).}
\footnotesize
\begin{tabular}{@{}lrrr@{}}
\toprule
\textbf{Method} & \textbf{Mean Cost} & \textbf{vs GT} & \textbf{vs Raw} \\
\midrule
DIFUSCO raw              & 6.696 & $+$13.67\% & \textit{baseline} \\
DIFUSCO $+$ 2-opt        & 5.969 & $+$1.32\%  & $-$10.87\% \\
\textbf{DIFUSCO$\rightarrow$DualOpt} & \textbf{5.793} & \textbf{$-$1.67\%} & \textbf{$-$13.50\%} \\
C$+$2opt (baseline)      & 5.891 & \textit{baseline} & $-$12.03\% \\
\bottomrule
\end{tabular}
\end{table}

The pipeline achieves \textbf{$-$1.67\% vs GT}---the best result of any method tested and the only one to consistently beat the training labels. The DualOpt reviser extracts 24\% more improvement than 2-opt from the same heatmap. On TSPLIB kroA100, the pipeline achieves 1.81\% gap, surpassing the original DualOpt (2.71\%) by 0.9~pp---demonstrating cross-distribution robustness. However, the pipeline \textbf{requires per-scale DIFUSCO training}: TSP-50 model on TSP-100 degrades ($-$2.24\% vs original DualOpt); TSP-100 model on TSP-100 recovers to $+$1.56\%.

\subsection{Five Improvement Strategies}

We systematically explored five strategies to integrate DIFUSCO's heatmap with DualOpt's reviser. Only the pipeline (\#2) succeeded; the other four failed, but each revealed fundamental design principles.

\vspace{2pt}
{\footnotesize
\begin{tabular}{@{}clll@{}}
\toprule
\textbf{\#} & \textbf{Strategy} & \textbf{Mechanism} & \textbf{Outcome} \\
\midrule
1 & Heatmap-Guided Reviser & Skip revision on confident windows & No effect ($\Delta\!=\!0.00\%$) \\
2 & \textbf{DIFUSCO$\rightarrow$DualOpt} & \textbf{Replace initial tour with diffusion} & \textbf{SUCCESS ($-$1.67\%)} \\
3 & Adaptive Window Sizing & 2-opt diagnostic $\rightarrow$ iteration budget & No improvement \\
4 & Fragment Freezing & Lock edges where two solvers agree & Mixed (4/10 improved, unstable) \\
5 & Destroy-and-Repair & Remove low-conf.\ edges, greedy reconnect & Degradation \\
\bottomrule
\end{tabular}
}
\vspace{4pt}

\textbf{Why \#1, \#3, \#4, \#5 Failed.} Three root causes: (1) \textit{Reviser saturation}---on TSP-50, DualOpt's reviser converges to a local optimum no perturbation can escape; (2) \textit{Heatmap miscalibration}---only $\sim$45\% of high-confidence edges are optimal even in-distribution; on TSPLIB: $\sim$0\%; (3) \textit{Weak signal}---2-opt (\#3) converges to the same local optimum as the reviser; solver consensus (\#4) agrees on wrong edges $\sim$40\% of the time.

\textbf{Engineering significance.} These five attempts systematically map the design space for integrating generative models with learned optimizers: (a) learned confidence $\neq$ actionable guidance (\#1); (b) changing the input beats constraining the optimizer (\#2 vs \#1,3,4,5); (c) a weaker optimizer cannot usefully guide a stronger one (\#3); (d) multi-solver structural consensus shows promise but needs a tiebreaker (\#4: 4/10 improved, up to $-$5.16\%); (e) once an optimizer saturates, perturbations only degrade (\#5).

% ======================================================================
\begin{thebibliography}{10}
\footnotesize

\bibitem{difusco} Z.~Sun and Y.~Yang. DIFUSCO: Graph-based Diffusion Solvers for Combinatorial Optimization. \textit{NeurIPS}, 2023.

\bibitem{dualopt} Z.~Zhou et al.\ DualOpt: A Dual Divide-and-Optimize Algorithm for the Large-Scale TSP. \textit{AAAI}, 2025.

\bibitem{cappart2023} Q.~Cappart, D.~Ch\'etelat, E.~Khalil, et al. Combinatorial optimization and reasoning with graph neural networks. \textit{JMLR}, 24(130):1--61, 2023.

\bibitem{gensco} M.~Kim, S.~Park, et al. GenSCO: Generative Sequence Combinatorial Optimization. \textit{NeurIPS}, 2024.

\end{thebibliography}

\end{document}
